%AssociationSemantics.tex

\ksec{Associations Semantics (KerML 8.4.4.5.1)}

\pn An \bemph{association} is a  {classifier} that classifies links between things.
At least two of association features must be association {ends}\index{end}.
An association is also a {relationship} between the types of its association ends, called its related types.
Associations can also have features that are not end features.

\pn An association must specialize, directly or indirectly, \texttt{Links::Link}.

\pn The end features of an association identify the participants in the links that are {instances}
of the association and must have multiplicity [1..1].

\pn Lexical ordering matters for associations.
For binary associations, the two related
types are identified as the \bemph{source} type and the \bemph{target} type (which may be the same).
For associations with more than
two association ends (``n-ary"), the first related type is the source type and all the remaining related types are target
types.

\ecaption{association}
\begin{lstlisting}
assoc Ownership { // Specializes Links::BinaryLink by default.
  end feature owner[1] : LegalEntity; // Redefines BinaryLink::source.
  end feature ownedAsset[1] : Asset; // Redefines BinaryLink::target.
  feature valuationOnPurchase : MonetaryValue;
}
\end{lstlisting}

\pn For a binary association, the source and target ends are already declared to have multiplicity [1..1], so this
does not need to be redeclared in redefinitions of these features. 

\begin{lstlisting}
assoc Ownership {
  end owner : LegalEntity;
  end ownedAsset : Asset;
  feature valuationOnPurchase : MonetaryValue;
}
\end{lstlisting}

\pn However, for non-binary associations, the end
multiplicity must be explicitly declared as [1..1] to override the usual default of [0..*].

\subsection{Binary Associations}

\pn Associations implicitly specialize \texttt{Links::Link}.

\ecaption{Link}
\begin{lstlisting}
    abstract assoc Link specializes Anything {
        doc
        /*
         * Link is the most general association between two or more things.
         */
        feature participant: Anything[2..*] nonunique ordered;
    }
\end{lstlisting}

\pn Abstract associations need not have ends, but nevertheless have at least two participants which may be anything.

\pn Associations with two ends implicitly specialize \texttt{Links::BinaryLink}.

\ecaption{BinaryLink}
\begin{lstlisting}
    assoc all BinaryLink specializes Link {
        doc
        /*
         * BinaryLink is the most general binary association between exactly two things, 
         * nominally directed from source to target.
         */       
        feature participant: Anything[2] nonunique ordered redefines Link::participant;   
        end feature source: Anything[1] subsets participant;
        end feature target: Anything[1] subsets participant;
    }
\end{lstlisting}

\pn Consider a binary association $A$ specializing association $L$ which must explicitly or implicitly specialize \texttt{BinaryLink} having two ends.
The first end (source) has identifier $s$ and type $S$.  The second end (target) has identifier $t$ and type $T$.  Association $A$ is represented as the class of triples having the first element \textbf{assoc}, the second element the wonce of ``A", and the third element the triple $\langle C,D,E\rangle$.
$C$ is the class for type $L$.  $D$ is the class having a single element $\langle\textbf{end},``s",\langle w,s \rangle\rangle$ for end feature $s$.
 $E$ is the class having a single element $\langle\textbf{end},``t",\langle w,t \rangle\rangle$ for end feature $t$.

\begin{definition}~\blTitle{df-binassoc}~   Define Binary Association
\begin{align}
\texttt{assoc}~ A~\texttt{:>}~L~\{ \texttt{end}~ s ~\texttt{:}~ S\texttt{;} \texttt{end}~ t ~\texttt{:}~ T\texttt{;} \}~\Rightarrow  \notag \\
\vdash \mc{A} = \{~\langle \textbf{assoc},w,\langle  \mc{C}, \mc{D}, \mc{E}\rangle\rangle ~\vert~
w=W(``A")~ \wedge 
 \mc{C}=\{~c~\vert ~c \in  \mc{L}~\}~\wedge~\notag \\
~s\in  \mc{S}~\wedge~ \mc{D}~=~\{~\langle\textbf{end},``s",\langle w,s \rangle\rangle~\}~\wedge \notag \\
~t\in  \mc{T}~\wedge~ \mc{E}~=~\{~\langle\textbf{end},``t",\langle w,t \rangle\rangle~\}~\}~\label{eq:df-binassoc}\tag{df-binassoc}
\end{align}
\end{definition}


\subsection{Cross Features}

\pn Cross features restrict associations so that linked classifiers must refer to each other.\index{cross feature}

\pn For a binary association, one or both association ends can be explicitly declared to subset across feature owned by
the other related type. 
This is done with cross subsetting, which is a special kind of subsetting relationship specified
using the keyword \texttt{crosses} or the symbol \texttt{=>}. 

\pn Only end features may have cross subsetting relationships, and an
end feature can have at most one owned cross subsetting.

\ecaption{ cross subsetting}
\begin{lstlisting}
classifier LegalEntity {
  feature assetsOwned [*] ordered : Asset;
}
classifier Asset {
  feature owningEntities [1..*] : LegalEntity;
}
assoc AssetOwnership {
  end owner : LegalEntity => ownedAsset.owningEntities;
  end ownedAsset : Asset => owner.assetsOwned;
  feature valuationOnPurchase : MonetaryValue;
}
\end{lstlisting}

\pn The cross subsetting by a binary association must be a feature chain starting with the other feature.
Thus the cross subsetting of \texttt{owner} is a feature chain starting with \texttt{ownedAsset}.
Conversely, the cross subsetting of \texttt{ownedAsset} is a feature chain starting with \texttt{owner}.

\pn Just because the cross feature restriction is met, does not mean an instance of the association exists.  The restriction is only necessary.

\pn Inserting the \texttt{all} keyword requires an association instance to exist whenever the cross feature restriction is met.  Then the restriction is sufficient.

\pn The definition of binary association cross feature is very similar to \ref{eq:df-binassoc} with the addition that the $f$ feature of $s$ equal $t$,
and the $g$ feature of $t$ equal $s$.

\begin{definition}~\blTitle{df-binassocx}~   Define Binary Association Cross Feature
\begin{align}
\texttt{classifier} ~S~ \{ \texttt{feature} ~f ~\texttt{:}~ T\texttt{;} \}~~~~ \notag \\
\texttt{classifier} ~T~ \{ \texttt{feature} ~g ~\texttt{:}~ S\texttt{;} \}~~~~ \notag \\
\texttt{assoc}~ A~\texttt{:>}~L~\{ \texttt{end}~ s ~\texttt{:}~ S~\texttt{=>}~t.g\texttt{;} \texttt{end}~ t ~\texttt{:}~ T~\texttt{=>}~s.f~\texttt{;} \}~\Rightarrow  \notag \\
\vdash  \mc{A} = \{~\langle \textbf{assoc},w,\langle  \mc{C},\mc{D},\mc{E}\rangle\rangle ~\vert~
w=W(``A")~ \wedge 
\mc{C}=\{~c~\vert ~c \in \mc{L}~\}~\wedge~\notag \\
~s\in \mc{S}~\wedge~s.f=t~\wedge~\mc{D}~=~\{~\langle\textbf{end},``s",\langle w,s \rangle\rangle~\}~\wedge \notag \\
~t\in \mc{T}~\wedge~t.g=s~\wedge~\mc{E}~=~\{~\langle\textbf{end},``t",\langle w,t \rangle\rangle~\}~\}~\label{eq:df-binassocx}\tag{df-binassocx}
\end{align}
\end{definition}

\subsection{Owned Cross Features}

\pn It is also possible to declare cross features directly in the declaration of the ends of an association, rather than nested
in the related types of the association. 
Such owned cross features are declared between the \texttt{end} and \texttt{feature}
keywords of the association end declarations.

\pn Using cross subsetting:
\ecaption{cross subsetting}
\begin{lstlisting}
classifier LegalEntity {
  feature assetsOwned [*] ordered : Asset;
}
classifier Asset {
  feature owningEntities [1..*] : LegalEntity;
}
assoc AssetOwnership {
  end owner : LegalEntity => ownedAsset.owningEntities;
  end ownedAsset : Asset => owner.assetsOwned;
  feature valuationOnPurchase : MonetaryValue;
}
\end{lstlisting}

\pn Using owned cross features:
\ecaption{owned cross features}
\begin{lstlisting}
assoc AssetOwnership {
  end owningEntities[1..*] feature owner : LegalEntity;
  end assetsOwned[*] ordered feature ownedAsset : Asset;
  feature valuationOnPurchase : MonetaryValue;
}
\end{lstlisting}

\pn In fact, the feature identifiers \texttt{owningEntities} and \texttt{assetsOwned} are unnecessary.  
Which feature crossed is determined by type and multiplicity.

\ecaption{AssetOwnership}
\begin{lstlisting}
assoc AssetOwnership {
  end [1..*] feature owner : LegalEntity;
  end [*] ordered feature ownedAsset : Asset;
  feature valuationOnPurchase : MonetaryValue;
}
\end{lstlisting}

\pn What would happen if \texttt{LegalEntity} could lease assets in addition to owning them?

\ecaption{LegalEntity}
\begin{lstlisting}
classifier LegalEntity {
  feature assetsOwned [*] ordered : Asset;
  feature assetsLeased [*] ordered : Asset;
}
\end{lstlisting}

\pn Binary associations having owned cross features have the same semantics as \ref{eq:df-binassocx}.


\subsection{Self-Link}

\pn Self link asserts its end are the same (identity).  Curiously, \texttt{SelfLink} uses both cross subsetting and owned cross feature.
Seems like \texttt{thisThing} and \texttt{sameThing} should be sufficient to assert identity.  What does \texttt{self2} mean?
Shouldn't it be just \texttt{self}?

\ecaption{SelfLink}
\begin{lstlisting}
 assoc all SelfLink specializes BinaryLink {
    doc
    /*
     * SelfLink is a binary association in which the things at the two ends are asserted
     * to be the same.
     */      
    language "Domain" /* thisThing = sameThing */
    end feature thisThing: Anything redefines source subsets sameThing crosses sameThing.self;
    end self2 [1] feature sameThing: Anything redefines target subsets thisThing;
}
\end{lstlisting}

\subsection{N-ary Associations}\label{sec:naryassoc}

\pn Associations having more than two ends (n-ary) specialize \texttt{Links::Link}.  Though the spec says the first end is the source, and subsequent ends are targets, there is nothing in \texttt{Links::Link} which names either.

\pn Formal semantics for n-ary associations generalize \ref{eq:df-binassoc}.

\begin{definition}~\blTitle{df-assoc}~   Define N-ary Association
\begin{align}
\texttt{assoc}~ A~\texttt{:>}~L~\{ \texttt{end}~ t_1 ~\texttt{:}~ T_1\texttt{;}~\ldots~ \texttt{end}~ t_n ~\texttt{:}~ T_n\texttt{;} \}~\Rightarrow  \notag \\
\vdash \mc{A} = \{~\langle \textbf{assoc},w,\langle \mc{C},\mc{D_1},\ldots,\mc{D_n}\rangle\rangle ~\vert~
w=W(``A")~ \wedge 
\mc{C}=\{~c~\vert ~c \in \mc{L}~\}~\wedge~\notag \\
~t_1\in \mc{T_1}~\wedge~\mc{D_1}~=~\{~\langle\textbf{end},``t_1",\langle w,t_1 \rangle\rangle~\}~\wedge \notag \\
~~~\ldots~\wedge~~~~~~~~~~~~~ \notag \\
~t_n\in \mc{T_n}~\wedge~\mc{D_n}~=~\{~\langle\textbf{end},``t_n",\langle w,t_n \rangle\rangle~\}~\}~\label{eq:df-assoc}\tag{df-assoc}
\end{align}
\end{definition}

\pn The \texttt{crosses} or \texttt{=>} notation for cross features does not work for n-ary associations, but the owned cross features does.
But what could owned cross features of n-ary associations mean?  The association example \texttt{ProductSelection\_N-ary.kerml} gives an example for N=3.

\ecaption{N-ary Cross Features}
\begin{lstlisting}
	assoc ProductSelection {
		end [0..1] feature cart: ShoppingCart[1];
		end [0..*] feature selectedProduct: Product[1];
		end [1..1] feature account : Account[1];
	}
\end{lstlisting}

\pn The example is elaborated with named end features, then nested cross features, and finally implied cross subsetting of Cartesian-product features.  Given end multiplicities other than [1] its wildly convoluted.  Don't bother with cross features of n-ary associations; you'll just make an incomprehensible mess.

\ecaption{AgreedAssetOwnership}
\begin{lstlisting}
assoc AgreedAssetOwnership {
  end [*] feature owner : LegalEntity;
  end [*] ordered feature ownedAsset : Asset;
  end [0..1] feature agreement : OwnershipAgreement;
}
\end{lstlisting}

\pn The cross multiplicity [0..1] for \texttt{agreement} requires that, for every pair of a \texttt{LegalEntity} and an \texttt{Asset}, at most
one \texttt{AgreedAssetOwnership} instance may link that \texttt{LegalEntity} as its \texttt{ownerandAsset} as its \texttt{ownedAsset} to
an \texttt{OwnershipAgreement} as its \texttt{agreement}. Similarly, every pair of a \texttt{LegalEntity} as \texttt{owner} and an
\texttt{OwnershipAgreement} as \texttt{agreement} may be linked to any number of \texttt{Assets} as \texttt{ownedAsset}, including none at
all, as declared by the cross multiplicity of \texttt{ownedAsset}. This collection of linked \texttt{Assests} is ordered, as declared
by the cross ordering of \texttt{ownedAsset}. The same applies to the cross multiplicity of \texttt{owner}.

\ksec{Association Structures (KerML 8.4.4.5.2)}

\texttt{AssociationStructure} is both \texttt{Association} and \texttt{Structure}, and specializes (directly, indirectly, or implicitly) \texttt{Objects::LinkObject} which specializes 
both \texttt{Links::Link} and \texttt{Objects::Object}.

Therefore an association structure instance is an \texttt{Occurrence} having duration.

Associations are timeless.  Association structure instances have existence.


